\documentclass[11pt]{article}
\usepackage[a4paper,margin=23mm]{geometry}
\usepackage[T1]{fontenc}
\usepackage{lmodern,amsmath,amssymb,amsthm,mathtools,microtype,booktabs}
\usepackage[colorlinks=true,linkcolor=blue,citecolor=blue,urlcolor=blue]{hyperref}
\newcommand{\E}{\mathbb E}
\newcommand{\tr}{\operatorname{tr}}
\newcommand{\Nc}{\mathcal N}
\newcommand{\Pin}{\mathcal P_H}
\newcommand{\norm}[1]{\left\lVert#1\right\rVert}
\newtheorem{theorem}{Theorem}
\newtheorem{lemma}{Lemma}
\title{A conserved comparison for fixed-charge memory returns\\
\large Force-defect inversion and strong finite-time bounds}
\author{Alan Day \and Chengheng\thanks{AI research assistant (OpenAI).}}
\date{8 November 2026\quad Local research draft}
\begin{document}\maketitle
\begin{abstract}
The ordinary charged-tail estimate and the neutral Jacobi expansion
give a direct fixed-sector treatment of the first memory return.
An exact inverse of the body-loss identity bounds the force-sandwich
defect at every fixed finite Schatten order. A finite-parameter energy
function then approximates the scaled ordinary kernel at rate
$O(p^{-1/2})$. The first return is a conserved background up to the
same error; all later returns contribute $O(p^{-1})$. Exact energy
pinching therefore reduces the total memory field to ordinary
propagation in every fixed finite Schatten norm, uniformly on a
compact central charge window and compact real-time intervals.
\end{abstract}
\section{Model and result}
Use standard CAR, $\mathcal H_m=\wedge^m\mathbb C^N$,
$d_m=\binom Nm$, and the common static Hermitian Gaussian model
\[
 H=\sum_{|I|=|J|=p}J_{IJ}c_I^\dagger c_J,\quad
 \E J_{IJ}J_{KL}=w_N\delta_{IL}\delta_{JK},\quad
 w_N=\frac{4^pJ_s^2}{\lambda_N}\left(\frac{p!}{N^p}\right)^2.
\]
Overlaps are retained. Let $\lambda_N=p^2/N\to\lambda>0$,
$J_s>0$ fixed, $\hbar=1$, and use string time. Put
\[
 \norm X_{r,m}=(\E\tr_m|X|^r)^{1/r},\qquad
 C_m=w_N\binom mp\binom{N-m+p}p.
\]
All bounds are uniform on $|m-N/2|\le L_0p$ for fixed $L_0$,
and on $|t|\le S$ for fixed $S$, with fixed $1\le r<\infty$.
The known fixed-charge setting is described in \cite{bnr}.
Write $h=H_m$ and define the exact dilation
\begin{gather*}
 V\psi=m^{-1/2}\sum_i|i\rangle c_i\psi,\quad
 K=I_N\otimes H_{m-1},\quad Q=I-VV^\dagger,\\
 a=V^\dagger KV=\alpha h,\quad\alpha=1-p/m,\quad
 R=QKV,\quad D_\perp=QKQ,\\
 M(t)=R^\dagger e^{itD_\perp}R,\qquad
 \Lambda(t)=R^\dagger e^{itK}R.
\end{gather*}
The energy projection $\Pin X=\sum_eP_eXP_e$ retains entire
degenerate blocks. Let
\[
 \mathcal F(t)=(1-\Pin)m e^{-iht/2}M(t)e^{-iht/2},\qquad
 X_0(t)=(1-\Pin)m e^{-iht/2}\Lambda(t)e^{-iht/2}.
\]
\begin{theorem}\label{main}
In the stated model, norms and limits,
$\sup_{m,t}\norm{\mathcal F(t)-X_0(t)}_{r,m}=O(p^{-1/2})$.
\end{theorem}

\newpage
\section{An exact inverse of the ordinary body-loss identity}
On full neutral block families let
\[
 (DX)_m=mX_m-\sum_i c_i^\dagger X_{m-1}c_i,
 \qquad K_{n,m}=[p\Nc-D]H^n\big|_m.
\]
Here $\Nc$ is Gaussian coupling number. Define the random force
map and deterministic covariance map with their different inputs:
\begin{align*}
 \Gamma_{H,m}(X_{m-1})
   &=\sum_i[H,c_i]^\dagger X_{m-1}[H,c_i],\\
 \Phi_m(Y_m)&=\E_{H'}H'_mY_mH'_m,\\
 F_{n,m}&=\Gamma_{H,m}(H_{m-1}^n)-2p\Phi_m(h^n).
\end{align*}
The independent $H'$ only defines $\Phi$; all actual fields
use the same original disorder. The established body product rule
and Gaussian Euler--Laplacian identity \cite{operator,gradient} give
\begin{equation}\label{forward}
 K_{n+2,m}=\sum_{u+v+w=n}h^u F_{v,m}h^w,
 \qquad K_{0,m}=K_{1,m}=0.
\end{equation}
Indeed, both $DH^n$ and $p\Nc H^n$ start with $pnH^n$;
their second-order terms are respectively $\Gamma_H$ and $2p\Phi$.
\begin{lemma}\label{inverse}
For every $n\ge0$ the exact inverse is
\begin{equation}\label{inverse-eq}
 F_{n,m}=K_{n+2,m}-hK_{n+1,m}-K_{n+1,m}h+hK_{n,m}h.
\end{equation}
Moreover $\norm{F_{n,m}}_{r,m}\le A_r(n+1)^{b_r}L_r^n\sqrt{n!}$
with constants independent of $N,m,n$ on the stated window.
\end{lemma}
\begin{proof}
In \eqref{forward}, subtracting $hK_{n+1}$ leaves only $u=0$.
Subtracting $K_{n+1}h-hK_nh$ then removes all $w\ge1$.
The argument includes $n=0,1$ by $K_0=K_1=0$ and never
commutes a force defect through $h$.

Equation (O) of \cite{gradient} bounds $K_n$ itself by a Gaussian
factorial envelope. Alternatively, its main theorem bounds
$K_n-\lambda_N\Nc(\Nc-1)H^n$, with $C_m^{-n/2}$ applied after
loss. The OU/Chebyshev multiplier bound \cite{operator} gives
$\norm{\Nc(\Nc-1)H^n}_{r,m}\le c n^2(1+\sqrt2)^n
\norm{H^n}_{r,m}$. Together with $C_m\le C_*$ and
$\norm{H^n}_{r,m}\le L_r^n\sqrt{n!}$ this gives the same bound
on $K_n$. In \eqref{inverse-eq}, use $2r$ Holder for each one-sided
product and $4r$ for $hK_nh$. The extra factor
$\sqrt{(n+2)!/n!}$ is polynomial. Only a fixed multiple of r is
needed, independent of n.
\end{proof}
Thus the entire exponential series is summable in these norms:
\begin{equation}\label{force-time}
 F_m(t)=\Gamma_{H,m}(e^{itH_{m-1}})-2p\Phi_m(e^{ith}),
 \qquad \sup_{m,|t|\le S}\norm{F_m(t)}_{r,m}\le A_{r,S}.
\end{equation}
The tail is bounded by $A_r\sum_{n>L}(n+1)^{b_r}
(L_rS)^n/\sqrt{n!}$. Finite-N identification with the displayed
maps follows first by the exponential Taylor series; Gaussian
moments control its products with the two linear force factors.

\newpage
\section{A finite-parameter conserved covariance comparison}
Retain the exact finite-N quantities from \cite{neutral}:
\begin{gather*}
 \theta_{j,m}=\prod_{u=0}^{p-1}
 \left(1-\frac{jp}{m-u}\right)_+
 \left(1-\frac{jp}{N-m+p-u}\right)_+,\\
 Q_{0,m}(x)=1,\quad Q_{1,m}(x)=x,\quad
 Q_{j+1,m}(x)=xQ_{j,m}(x)-b_{j,m}Q_{j-1,m}(x),\\
 b_{j,m}=C_m\sum_{v<j}\theta_{v,m},\qquad
 T_{j,m}=p(\Phi_m-C_m\theta_{j,m})Q_{j,m}(h).
\end{gather*}
On the central window $b_{j,m}$ are uniformly bounded. The
exact entire expansion and saved bounds are
\begin{align}
 e^{ith}&=\sum_{j\ge0}c_{j,m}(t)Q_{j,m}(h),\nonumber\\
 |c_{j,m}(t)|&\le A_S L_S^j/j!,\qquad
 \norm{T_{j,m}}_{r,m}\le A_r(j+1)^{b_r}L_r^j\sqrt{j!}.
 \label{jacobi-bounds}
\end{align}
The polynomials $Q_j(h)$ have an analogous factorial envelope,
by their bounded recurrence coefficients and Gaussian moments.
Consequently each of the following series converges in every
fixed finite norm, uniformly on compact real times:
\begin{equation}\label{fdef}
 f_m(t,h)=\sum_{j\ge0}c_{j,m}(t)\theta_{j,m}Q_{j,m}(h),
 \qquad Z_m(t)=\sum_{j\ge0}c_{j,m}(t)T_{j,m}.
\end{equation}
The exact identity before pinching is
\begin{equation}\label{phif}
 \Phi_m(e^{ith})=C_m f_m(t,h)+p^{-1}Z_m(t).
\end{equation}
For each sample the scalar Jacobi series is entire in h, so
$f_m(t,h)$ is a function of h, including at degenerate eigenvalues.
Also $f_m(-t,h)=f_m(t,h)^\dagger$.

The map $\Phi_m/C_m$ is unital, completely positive and trace
preserving. It contracts Schatten norms, and maps a unitary to
an operator of norm at most one. Since $C_m$ is bounded away
from zero by $C_->0$ on the window, \eqref{phif} implies
\begin{equation}\label{fbnd}
 \sup_{m,|t|\le S}\norm{f_m(t,h)}_{r,m}
 \le 1+\frac{A_{r,S}}{pC_-}=O_{r,S}(1).
\end{equation}
No samplewise operator bound on f is required. Define
\begin{equation}\label{bdef}
 \beta_m=\frac{2p^2}{m}C_m,\qquad
 b_m(t,h)=\beta_m f_m(t,h).
\end{equation}
Then $[b_m(t,h),h]=0$ exactly and all its fixed finite norms
are uniformly bounded, since $p^2/m=O(1)$.

This comparison uses finite coefficients rather than replacing
them by $q^j$. The standard q-Gaussian Mehler comparison used
in \cite{background,bks} remains compatible with the limit,
but is not an additional input to the strong estimate here.

\newpage
\section{Ordinary propagation: every fixed finite norm}
Let $F_i=[H,c_i]$, $\mathcal A_m=\sum_iF_i^\dagger F_i|_m$,
$D_F=KV-Vh$, and $G=R^\dagger R$. Exact identities give
\begin{align}
 D_F^\dagger D_F&=\mathcal A_m/m,&
 R&=D_F+\delta Vh,\quad\delta=p/m,\nonumber\\
 D_F^\dagger e^{itK}D_F&=\Gamma_{H,m}(e^{itH_{m-1}})/m,&
 0\le G&=\mathcal A_m/m-\delta^2h^2\le\mathcal A_m/m.
 \label{dilation}
\end{align}
All rectangular norms use the same source dimension $d_m$,
even on the larger dilation space. Thus there is no change of
trace normalization when applying rectangular Holder.

The sectorwise Gaussian force moments of \cite{return} are
\begin{equation}\label{force-mom}
 \E\tr_m\mathcal A_m^k\le(2k-1)!!(Nv_N)^k,
 \quad v_N\le\frac{4J_s^2}{\lambda_N}\frac pN e^{2\lambda_N}.
\end{equation}
They follow from both force covariance bounds, Hermitian
dilation, Gaussian integration by parts, and Minkowski over i;
the force components need not be independent. Hence, using a
larger fixed integer moment when necessary,
\begin{gather}
 \norm{\mathcal A_m}_{r,m}=O_r(p),\quad
 \norm G_{r,m}=O_r(p/m)=O_r(p^{-1}),\label{Gbnd}\\
 \norm{D_F}_{2r,m}=O_r(p^{-1/2}),\qquad
 \norm{\delta Vh}_{2r,m}=O_r(p^{-1}).\label{DFbnd}
\end{gather}
The G estimate uses eigenvalue monotonicity, not operator
monotonicity of positive powers. Unitary propagation yields
$\norm{\Lambda(t)}_{r,m}\le\norm G_{r,m}$.

Expanding $R=D_F+\delta Vh$ and using \eqref{DFbnd} gives
\begin{align}
 &\norm{p\Lambda(t)-pD_F^\dagger e^{itK}D_F}_{r,m}\nonumber\\
 &\qquad\le p\bigl(2\norm{D_F}_{2r,m}\norm{\delta Vh}_{2r,m}
                         +\norm{\delta Vh}_{2r,m}^2\bigr)
 =O_r(p^{-1/2}).\label{shift-error}
\end{align}
Combine \eqref{force-time}, \eqref{phif}, \eqref{bdef} and
\eqref{dilation} without changing the lower-charge input:
\begin{equation}\label{force-b}
 \frac p m\Gamma_{H,m}(e^{itH_{m-1}})-b_m(t,h)
   =\frac p m F_m(t)+\frac{2p}{m}Z_m(t).
\end{equation}
Both terms are $O_{r,S}(p^{-1})$, because $m\asymp p^2$.
It follows that
\begin{equation}\label{leading}
 \sup_{m,|t|\le S}\norm{p\Lambda(t)-b_m(t,h)}_{r,m}
       =O_{r,S,L_0}(p^{-1/2}).
\end{equation}
The rate holds at each fixed finite r directly; it is not an
interpolation of an $L^2$ rate with an $O(p)$ coarse bound.
It concerns a leading ordinary kernel. Multiplying
\eqref{leading} by $m/p$ would not give a bounded dynamic field.

\newpage
\section{The first return and its exact remainder}
Write $U(t)=e^{iat}$ and use time convolution preserving matrix
order. The finite-dimensional block Duhamel identity from
\cite{return} is
\[
 \Lambda=M-M*U*\Lambda,\qquad
 M=\Lambda+J_1+E_1,
\]
where $J_1=\Lambda*U*\Lambda$ and
$E_1=M*U*\Lambda*U*\Lambda$. The minus sign in the first
identity comes from two transitions and $i^2=-1$.
Auxiliary-kernel removal is an established Mori method
\cite{mori}; here only finite matrices and finite iteration are used.

For $s\ge0$, let $\Delta_s=\{v,w\ge0:v+w\le s\}$ and
$u=s-v-w$. Then
\begin{equation}\label{first}
 J_1(s)=\int_{\Delta_s}\Lambda(v)U(u)\Lambda(w)\,dv\,dw.
\end{equation}
Define the finite-parameter conserved reference
\begin{equation}\label{kdef}
 \mathcal K_m(s,h)=\frac m{p^2}e^{-ihs}
 \int_{\Delta_s}b_m(v,h)e^{iau}b_m(w,h)\,dv\,dw.
\end{equation}
Every factor commutes with h. Each fixed finite norm of
$\mathcal K_m$ is uniformly bounded by Holder and $m/p^2=O(1)$.
Subtract the integrands in \eqref{first} and \eqref{kdef} as
\[
 (p\Lambda(v)-b_m(v,h))U(u)(p\Lambda(w))
       +b_m(v,h)U(u)(p\Lambda(w)-b_m(w,h)).
\]
Apply Holder at $2r$, \eqref{leading}, and the bounds on
$p\Lambda,b_m$. The simplex has area $s^2/2$, so
\begin{equation}\label{first-error}
 \sup_{m,0\le s\le S}
 \norm{m e^{-ihs/2}J_1(s)e^{-ihs/2}-\mathcal K_m(s,h)}_{r,m}
       =O(p^{-1/2}).
\end{equation}
No operator-norm bound on the comparison function is used.

The integrand of $E_1$ contains six factors R or $R^\dagger$
and unitary propagators. Holder at $6r$ on the common source
normalization, including Gaussian expectation, gives
\begin{equation}\label{E1}
 \norm{E_1(s)}_{r,m}\le\frac{|s|^4}{4!}\norm G_{3r,m}^3,
 \qquad \sup_{m,|s|\le S}\norm{mE_1(s)}_{r,m}=O(p^{-1}).
\end{equation}
This is an exact finite-return remainder, not a time Taylor
truncation. Reverse the simplex variables to get
$J_1(-s)=J_1(s)^\dagger$ and define
$\mathcal K_m(-s,h)=\mathcal K_m(s,h)^\dagger$.
The same bounds hold on both time signs. Combining them yields
\begin{equation}\label{unpinched}
 \sup_{m,|s|\le S}
 \norm{m e^{-ihs/2}(M(s)-\Lambda(s))e^{-ihs/2}
                         -\mathcal K_m(s,h)}_{r,m}=O(p^{-1/2}).
\end{equation}
Finally $(1-\Pin)\mathcal K_m=0$ samplewise and
$\norm{1-\Pin}_{r\to r}\le2$. This proves Theorem~\ref{main}.

\newpage
\section{Preparations, calibration and the next interface}
At half filling let $P_O$ be any already controlled fixed
Gaussian-degree selected polynomial preparation. Its finite-N
orbit density from \cite{density} satisfies
\[
 \E\langle P_O^\dagger YP_O\rangle_O=E\tr_m(R_{P,N}Y),
 \qquad\sup_N\norm{R_{P,N}}_{2,m}<\infty
\]
for covariant integrable Y. Set $\mathcal E=\mathcal F-X_0$.
It is covariant, including its exact energy pinching. Thus
\begin{equation}\label{prep}
 \E\norm{\mathcal E(s)P_O|O\rangle}^2
 \le\norm{R_{P,N}}_{2,m}\norm{\mathcal E(s)}_{4,m}^2
 =O(p^{-1}).
\end{equation}
This is a strong return-error estimate on the fixed polynomial
preparation core. It does not require interchanging exact
pinching and the large-N limit. It also does not by itself
bound arbitrary products of actual total-memory fields: their
other factors still need uniform high moments.

The reference is nonzero before pinching. Since
$f_m(0,h)=1$ and $\partial_tf_m(0,h)=i\theta_{1,m}h$,
\eqref{kdef} gives, at fixed N,
\begin{align}
 \mathcal K_m(s,h)
  ={}&\frac{2p^2}{m}C_m^2s^2\,I\nonumber\\
    &+i\frac{2p^2}{3m}C_m^2
           (2\theta_{1,m}+\alpha-3)hs^3+O(s^4).
 \label{calibration}
\end{align}
To obtain the cubic term use
$\int_{\Delta_s}v=\int_{\Delta_s}w=\int_{\Delta_s}u=s^3/6$,
and retain the outside $e^{-ihs}$. The fixed-order coefficients
have bounded finite moments. At half filling,
$p^2/m\to2\lambda$, $C_m\to C=J_s^2/\lambda$,
$\theta_{1,m}\to q=e^{-4\lambda}$, and $\alpha\to1$.
The first two coefficients therefore reproduce the existing
background calibration of \cite{background}:
\[
 4\lambda C^2s^2I+i\frac{8\lambda C^2}{3}(q-1)hs^3.
\]
This consistency check is not a new finite-time Mehler formula.

The proved reduction is $\mathcal F-X_0=O(p^{-1/2})$ in every
fixed finite norm, in the actual fixed charge sectors. The next
task is the fluctuation-scale two-component description of
$X_0$. It still requires the central four-auxiliary-field
conditional regression and the charged finite-$\alpha$ interface.
The actual bilocal Volterra reconstruction must then be
transferred with its required high moments. Arbitrary interlaced
nonpolynomial selected blocks, the spatial/mass dictionary and
finite-G gravitational dynamics are not conclusions of this note.

\newpage
\section{Saved finite checks and attribution}
The independent finite checks retain all Gaussian chaoses and
both adjacent charge sectors. They do not replace the analytic
bounds or establish a double-scaling limit by numerical fitting.
\begin{center}
\begin{tabular}{lll}\toprule
Check&Parameters&Coverage\\\midrule
Free noncommutative inverse&$n=0,\ldots,20$&21 exact integer rows\\
CAR/Hermite inverse&$(N,p,m)=(3,1,2)$&$n=0,1,2,3$\\
&$(4,1,2)$&$n=0,1,2$\\
&$(4,2,3)$&$n=0,1$\\
Block propagation&same three systems&five real times each\\
Return norm inequalities&$r=1,2,4$&90 inequalities\\\bottomrule
\end{tabular}
\end{center}
All nine CAR rows have a nonzero lower-chaos part. The real
Gaussian coordinate basis is normalized by $C_m=1$, with the
same coordinates in charge m and m-1. The block checks use
seeds 731, 732, 733 and times $-1.1,-0.2,0,0.4,1.1$.
Upper-triangular block matrix exponentials evaluate the ordered
convolutions directly, with no quadrature or time cutoff.
They check $M=\Lambda+J_1+E_1$, the force-dilation identity,
and the source-normalized bounds on $J_1,E_1$.
Floating residuals and tolerances are retained in the JSON;
they are not interval certificates. The formal free-word
cancellation uses exact integer arithmetic.

The body-loss and Gaussian product identities, the ordinary
factorial tail, the neutral Jacobi decomposition and the finite
Volterra expansion are inherited results. This note combines
them through the exact inverse and a finite-parameter conserved
comparison to obtain the sectorwise strong return estimate.
No priority claim is made for the elementary inverse or the
standard operator inequalities. The local evidence record binds
the code, scientific output, proof inputs, compilation and pages.

\begin{thebibliography}{9}\small
\bibitem{gradient} A. Day and Chengheng, \emph{Fixed-charge gradients
and the ordinary charged tail}, local draft, 8 November 2026;
\texttt{manuscript/susskind-sector-gradient-tail/main.tex}, (O), (T).
\bibitem{neutral} A. Day and Chengheng, \emph{The neutral memory field
in a fixed charge sector}, local draft, 8 November 2026;
\texttt{manuscript/susskind-sector-neutral-transfer/main.tex}, Sec.~4.
\bibitem{operator} A. Day and Chengheng, ordinary fourth-norm tail
operator draft, 7 November 2026;
\texttt{notes/susskind-bilocal-fourth-tail-operator-draft.md}, (2), (21).
\bibitem{return} A. Day and Chengheng, \emph{Finite memory times:
a controlled two-return reduction}, local draft, 6 November 2026;
\texttt{manuscript/susskind-memory-return-expansion/main.tex}, Secs.~2--3.
\bibitem{background} A. Day and Chengheng, \emph{The conserved
background of the first memory return}, local draft, 6 November 2026;
\texttt{manuscript/susskind-memory-return-background/main.tex}.
\bibitem{density} A. Day and Chengheng, fixed polynomial preparation
orbit-density draft, 8 November 2026;
\texttt{manuscript/susskind-slater-polynomial-density/main.tex}.
\bibitem{bnr} M. Berkooz, V. Narovlansky and H. Raj,
\emph{Complex Sachdev-Ye-Kitaev model in the double scaling limit},
\href{https://arxiv.org/abs/2006.13983}{arXiv:2006.13983}, Secs.~5, 7.
\bibitem{bks} M. Bo\.{z}ejko, B. K\"ummerer and R. Speicher,
\emph{q-Gaussian processes: non-commutative and classical aspects},
\href{https://arxiv.org/abs/funct-an/9604010}{arXiv:funct-an/9604010}.
\bibitem{mori} A. Montoya-Castillo and D. R. Reichman,
\emph{Approximate but Accurate Quantum Dynamics from the Mori
Formalism: I. Nonequilibrium Dynamics},
\href{https://arxiv.org/abs/1603.01903}{arXiv:1603.01903}, Sec.~III.
\end{thebibliography}
\end{document}
